\chapter{Feasibility Study}

\section{Gantt Chart and Task Definition}

The project comprises \textbf{37 tasks} across 7 categories, totalling 820 hours
(20 ECTS $\times$ 25 h/ECTS). The detailed phase breakdown and activity descriptions
are provided in the Methodology chapter (Chapter~4).

\begin{table}[H]
    \centering
    \begin{tabularx}{\textwidth}{Xrr}
\toprule
    \textbf{Category} & \textbf{Tasks} & \textbf{Hours} \\
\midrule
    Preparation      & 4  & 105h \\
    Analysis         & 5  & 49h  \\
    Documentation    & 6  & 95h  \\
    Development      & 8  & 178h \\
    Testing          & 5  & 123h \\
    Implementation   & 6  & 175h \\
    Milestones       & 3  & 95h  \\
    \textbf{TOTAL}   & \textbf{37} & \textbf{820h} \\
\bottomrule
    \end{tabularx}
    \caption{Distribution of tasks by category}
    \label{tab:task_categories}
\end{table}

The complete Gantt chart was built with a custom interactive web tool available at
\path{docs/resources/gantt/index.html} in the project repository.

\subsection{Critical Path}

The project has \textbf{13 critical tasks} (277 h total) that directly affect milestone deadlines:

\begin{table}[H]
  \centering
  \small
  \setlength{\tabcolsep}{4pt}
  \renewcommand{\arraystretch}{1.1}
  \begin{tabularx}{\textwidth}{lXcc}
    \toprule
    \textbf{ID} & \textbf{Critical Task} & \textbf{Hours} & \textbf{Deadline} \\
    \midrule
    1.4   & Pentesting Methodology and Validation           & 8h  & Nov 26, 2025 \\
    1.10  & Milestone: Preliminary Project                  & 40h & Jan 16, 2026 \\
    2.1   & Lab 1 Research                                  & 14h & Jan 26, 2026 \\
    2.2   & Lab 2 Research                                  & 14h & Jan 31, 2026 \\
    2.3   & Lab 3 Research -- Network Analysis \& IR        & 14h & Feb 7, 2026  \\
    2.6   & Pre-validation with Pilot Users                 & 20h & Feb 23, 2026 \\
    2.6.1 & Pre-validation Revision                         & 6h  & Mar 1, 2026  \\
    3.1   & Lab 4 Research                                  & 16h & Mar 23, 2026 \\
    3.4   & Penetration Testing Round 1                     & 30h & Apr 16, 2026 \\
    3.5   & Penetration Testing Round 2                     & 25h & Apr 30, 2026 \\
    3.9   & Bachelor's thesis Report Writing I              & 15h & May 1, 2026  \\
    3.10  & Bachelor's thesis Report Writing II             & 40h & May 20, 2026 \\
    3.11  & Bachelor's thesis Report Writing III            & 35h & May 25, 2026 \\
    \midrule
    \multicolumn{2}{l}{\textbf{TOTAL CRITICAL PATH}} & \textbf{277h} & \\
    \bottomrule
  \end{tabularx}
  \caption[Critical path tasks]{Critical path tasks and deadlines}
  \label{tab:critical_task}
\end{table}

\section{Budget}

\subsection{Human Resources}

\begin{table}[H]
    \centering
    \begin{tabularx}{\textwidth}{Xrrr}
\toprule
    \textbf{Phase} & \textbf{Hours} & \textbf{Rate} & \textbf{Total} \\
\midrule
    Phase 0: Preparation   & 105h & 15\,\euro/h & 1{,}575\,\euro \\
    Phase 1: Preliminary Project & 188h & 15\,\euro/h & 2{,}820\,\euro \\
    Phase 2: Interim       & 212h & 15\,\euro/h & 3{,}180\,\euro \\
    Phase 3: Final         & 315h & 15\,\euro/h & 4{,}725\,\euro \\
    \textbf{TOTAL}         & \textbf{820h} & & \textbf{12{,}300\,\euro} \\
\bottomrule
    \end{tabularx}
    \caption{Human resources cost (junior engineer rate, 15\,\euro/h, theoretical)}
    \label{tab:human_resources_cost}
\end{table}

\subsection{Hardware and Software}

\begin{table}[H]
    \centering
    \begin{tabularx}{\textwidth}{Xrrr}
\toprule
    \textbf{Resource} & \textbf{Qty} & \textbf{Unit} & \textbf{Total} \\
\midrule
    HomeLab Server (2$\times$ Xeon E5-2680v4, 64\,GB, HDD+SSD) & 1 & 800\,\euro & 800\,\euro \\
    Development Workstation & 1 & 1{,}200\,\euro & 1{,}200\,\euro \\
    Network Equipment & 1 set & 150\,\euro & 150\,\euro \\
    External Storage & 1 & 100\,\euro & 100\,\euro \\
    \textbf{SUBTOTAL Hardware} & & & \textbf{2{,}250\,\euro} \\
\bottomrule
    \end{tabularx}
    \caption{Hardware costs}
    \label{tab:hardware_costs}
\end{table}

\begin{table}[H]
    \centering
    \begin{tabularx}{\textwidth}{Xrr}
\toprule
    \textbf{Resource} & \textbf{License} & \textbf{Cost} \\
\midrule
    EVE-NG Community Edition  & Free/Open Source & 0\,\euro \\
    Parrot OS Security        & Free/Open Source & 0\,\euro \\
    Metasploit Framework      & Free/Open Source & 0\,\euro \\
    Burp Suite Community      & Free             & 0\,\euro \\
    OWASP ZAP                 & Free/Open Source & 0\,\euro \\
    Wireshark                 & Free/Open Source & 0\,\euro \\
    Visual Studio Code        & Free             & 0\,\euro \\
    LaTeX (TeX Live)          & Free/Open Source & 0\,\euro \\
    Git + GitHub              & Free (Student)   & 0\,\euro \\
    \textbf{SUBTOTAL Software} & & \textbf{0\,\euro} \\
\bottomrule
    \end{tabularx}
    \caption{Software and license costs (all free/open-source for educational use)}
    \label{tab:software_costs}
\end{table}

The server is a one-time investment reused across multiple projects. An additional 32\,GB RAM
module (\textasciitilde{}80\,\euro) was acquired during the project and is not reflected
in the original estimate.

\subsection{Other Costs}

\begin{table}[H]
    \centering
    \begin{tabularx}{\textwidth}{Xr}
\toprule
    \textbf{Concept} & \textbf{Cost} \\
\midrule
    Electricity consumption (820\,h $\times$ 0.15\,kWh $\times$ 0.20\,\euro/kWh) & 25\,\euro \\
    Internet connectivity (9 months $\times$ 40\,\euro/month)                     & 360\,\euro \\
    Documentation printing and binding (3 copies)                                  & 60\,\euro \\
    Contingency fund (5\% of total)                                                & 635\,\euro \\
    \textbf{TOTAL Other Costs} & \textbf{1{,}080\,\euro} \\
\bottomrule
    \end{tabularx}
    \caption{Additional operational costs}
    \label{tab:other_costs}
\end{table}

\subsection{Total Budget}

\begin{table}[H]
    \centering
    \begin{tabularx}{\textwidth}{Xr}
\toprule
    \textbf{Category} & \textbf{Total} \\
\midrule
    Human Resources  & 12{,}300\,\euro \\
    Hardware         & 2{,}250\,\euro  \\
    Software         & 0\,\euro        \\
    Other (electricity, internet, contingency) & 1{,}080\,\euro \\
    \textbf{TOTAL}   & \textbf{15{,}630\,\euro} \\
\bottomrule
    \end{tabularx}
    \caption{Complete project budget}
    \label{tab:total_budget}
\end{table}

\section{Feasibility Analysis}

\subsection{Technical Feasibility}

Phase 0 infrastructure is 100\% operational (EVE-NG, Proxmox\cite{proxmox}, Parrot OS\cite{parrot_os},
development environment). All software is open-source and widely adopted in professional
cybersecurity training. The required competencies are:

\begin{table}[H]
    \centering
    \begin{tabularx}{\textwidth}{XXc}
\toprule
    \textbf{Area} & \textbf{Required Competency} & \textbf{Level} \\
\midrule
    Network Administration & VLANs, routing, firewall rules & Intermediate \\
    Operating Systems      & GNU/Linux administration (Debian/Ubuntu) & Intermediate \\
    Virtualisation         & EVE-NG topology design, VM management & Advanced \\
    Pentesting             & Reconnaissance, exploitation, web vulns & Intermediate \\
    Scripting              & Python and Bash for automation & Intermediate \\
    Documentation          & LaTeX, technical writing & Intermediate \\
\bottomrule
    \end{tabularx}
    \caption{Required technical competencies}
    \label{tab:technical_competencies}
\end{table}

Risk assessment follows NIST SP 800-30\cite{nist_sp80030}. Table~\ref{tab:risk_matrix}
summarises the six identified risks with probability, impact, and mitigation strategy.

\begin{table}[H]
  \centering
  \footnotesize
  \setlength{\tabcolsep}{3pt}
  \renewcommand{\arraystretch}{1.1}
  \begin{tabularx}{\textwidth}{XcccX}
    \toprule
    \textbf{Risk} & \textbf{Prob.} & \textbf{Impact} & \textbf{Level} & \textbf{Mitigation} \\
    \midrule
    EVE-NG Performance Issues       & Med & High & HIGH   & Dedicated 32\,GB RAM; VirtualBox fallback \\
    VM Image Incompatibility        & Med & Med  & MEDIUM & Early validation; QCOW2 standardisation \\
    Knowledge Gaps (advanced net.)  & Med & Med  & MEDIUM & Structured learning; tutor consultation \\
    Hardware Degradation            & Low & High & MEDIUM & Quarterly diagnostics; S.M.A.R.T. alerts \\
    Data Loss (backup insufficient) & Low & High & MEDIUM & Daily backups; Git + off-site cloud \\
    Configuration Gaps              & Med & Med  & MEDIUM & Continuous docs; topology exports via Git \\
    \bottomrule
  \end{tabularx}
  \caption[Risk assessment]{Risk assessment matrix (NIST SP 800-30). No risk is project-blocking.}
  \label{tab:risk_matrix}
\end{table}

\textbf{Conclusion}: High technical feasibility. All risks have concrete mitigations and none
is project-blocking.

\subsection{Economic Feasibility}

\begin{table}[H]
  \centering
  \small
  \setlength{\tabcolsep}{4pt}
  \renewcommand{\arraystretch}{1.1}
  \begin{tabularx}{\textwidth}{lXr}
    \toprule
    \textbf{Alternative} & \textbf{Description} & \textbf{Annual Cost} \\
    \midrule
    \textbf{This Project} & Custom EVE-NG labs, curriculum-aligned & \textbf{0\,\euro/year*} \\
    HackTheBox Academy    & Commercial labs platform               & 1{,}200\,\euro/year \\
    TryHackMe Education   & Gamified learning paths               & 800\,\euro/year     \\
    Cybrary for Teams     & Video training + labs                 & 1{,}500\,\euro/year \\
    Custom Enterprise Labs& Outsourced development                & 8{,}000--15{,}000\,\euro \\
    \bottomrule
  \end{tabularx}
  \caption[Cost comparison]{Cost comparison with commercial alternatives (*excludes initial development)}
  \label{tab:cost_comparison}
\end{table}

\textbf{Conclusion}: High economic feasibility. Zero recurring software costs; hardware
amortised over 3--5 years; comparable commercial alternatives cost 800--1{,}500\,\euro/year.


\subsection{Temporal Feasibility}

The project spanned 10 months (July 2025--May 2026) across four phases totalling
820 hours (20 ECTS $\times$ 25~h/ECTS), as shown in Table~\ref{tab:task_categories}.
All milestones were met on schedule: Phase~0 (infrastructure), Phase~1
(preliminary proposal), Phase~2 (Labs~1--3 implementation and interim
validation), and Phase~3 (Lab~4, final validation, and report writing).
The 820-hour workload was distributed as planned, with the critical path of
13 tasks (277~h, Table~\ref{tab:critical_task}) completed without blocking delays.

\textbf{Conclusion}: High temporal feasibility. The project was delivered
within the planned timeline and within the 820~h budget for 20 ECTS.

\subsection{Environmental Feasibility}

The project's energy footprint comes mainly from the HomeLab server (150\,W avg.) and
workstation (65\,W avg.). Development phase consumption: 188.6\,kWh; operational phase
(4 sessions $\times$ 30 students $\times$ 3\,h/year): 82.8\,kWh/year. A 5\,kWp rooftop
solar installation at the author's home generated 6{,}634.9\,kWh in 2025 (real data),
offsetting approximately 34\% of development-phase consumption and reducing CO\textsubscript{2}
emissions by 32\% compared to grid-only scenarios\cite{red_electrica,murugesan2008}.

\begin{table}[H]
  \centering
  \small
  \setlength{\tabcolsep}{4pt}
  \renewcommand{\arraystretch}{1.1}
  \begin{tabularx}{\textwidth}{lXrr}
    \toprule
    \textbf{Scenario} & \textbf{Description} & \textbf{CO\textsubscript{2}/year} & \textbf{With Solar} \\
    \midrule
    This Project (ops)  & Virtualised shared HomeLab           & 20.7\,kg & 14.1\,kg \\
    This Project (dev)  & Development phase, solar-offset      & 47.2\,kg & 32.0\,kg \\
    Individual VMs      & 30 students $\times$ personal laptop & 180\,kg  & --       \\
    Cloud-Based Labs    & AWS/Azure (PUE 1.5)                  & 95\,kg   & --       \\
    Physical Lab        & Dedicated hardware switches/routers  & 450\,kg  & --       \\
    \bottomrule
  \end{tabularx}
  \caption[CO\textsuperscript{2} comparison]{CO\textsuperscript{2} comparison across laboratory approaches}
  \label{tab:co2_comparison}
\end{table}

Sustainability measures: VMs suspended when idle; digital-only documentation (3 printed
copies required by regulation); modular server hardware designed for 3--5 year longevity;
open-source stack extends the usable lifespan of existing hardware.

\textbf{Conclusion}: Good environmental feasibility; significantly lower carbon footprint
than physical labs or individual student deployments.

\subsection{Legal Aspects}

\subsubsection{Intellectual Property and Licensing}

All software (EVE-NG, Parrot OS Security, Metasploitable, DVWA) is licensed under
permissive open-source licenses (GPL, MIT, Apache 2.0) that permit educational use,
modification, and redistribution. Per UPF regulations (Article 10), intellectual property
of the bachelor's thesis belongs to the student; exploitation rights default to the
student unless otherwise agreed with the tutor and institution. All citations follow
IEEE standards.

\subsubsection{Ethical and Legal Considerations for Pentesting}

All pentesting activities are conducted within isolated EVE-NG sandboxes with no
connection to production networks. Only intentionally vulnerable systems (Metasploitable,
DVWA\cite{dvwa}) designed for training are used as targets. The project is
conducted under academic supervision (Pere Vidiella i Catalan) within the scope of the
``Introduction to Cybersecurity'' curriculum. No personally identifiable information
is processed; all scenarios use synthetic data.

\subsubsection{Use of Artificial Intelligence Tools}

Auxiliary AI tools --- both locally deployed language models on the HomeLab and
complementary online services --- were used only to review draft text, detect
inconsistencies, and check documentation errors. No AI tool was used to generate
original technical content, conclusions, research findings, or design decisions;
all creative, analytical, and engineering work is the author's own. This disclosure
is provided in accordance with the academic integrity guidelines applicable to the
programme.

\subsubsection{Compliance with Academic Regulations}

The project follows Tecnocampus ESUPT bachelor's thesis regulations (approved
September 30, 2021): original work with IEEE citation standards (Article 11 on
plagiarism); structure aligned with Chapter 9 evaluation requirements (continuous
assessment: preliminary project proposal, interim, final delivery, tribunal defence); formatting
per the ESUPT redaction guide.

\subsection{Gender and Diversity Perspective}

\subsubsection{Inclusive Language and Representation}

All documentation --- thesis, technical write-ups, and student guides --- uses
gender-neutral language. Fictional accounts, scenario names, and user credentials
in lab environments are culturally neutral and carry no gender or demographic
assumptions. Lab scenarios use diverse names that reflect the multicultural nature
of the cybersecurity profession.

\subsubsection{Addressing the Gender Gap in Cybersecurity}

Cybersecurity remains a male-dominated field, with women representing around 25\%
of the global workforce\cite{isc2_workforce}. This project contributes to lowering
barriers to entry in two ways. First, well-structured and clearly documented labs
reduce the intimidation factor often associated with hands-on security tools,
making the material more approachable for students from diverse backgrounds.
Second, the ethical framing of the course --- emphasising responsible disclosure,
defensive security, and societal impact --- broadens the appeal of cybersecurity
beyond a narrow stereotype, which research suggests helps attract a more diverse
student population\cite{conti2019hands}.

\subsubsection{Accessibility Measures}

The practical barriers to engagement in this course are not demographic but relate
to two specific factors:

\begin{itemize}
    \item \textbf{Prior knowledge}: Students with limited background in networking or
          operating systems may find early labs more demanding. This is mitigated through
          structured guides with progressive hints and a tools reference section
          in each student PDF.

    \item \textbf{Reading difficulties (dyslexia)}: Dense technical documentation
          can be harder to process for dyslexic readers. This thesis is typeset in
          OpenDyslexic; lab documentation uses short paragraphs, descriptive headings,
          and visually separated code blocks. The pilot study confirmed that a participant
          with documented dyslexia was able to complete Lab~3 with the materials provided.
\end{itemize}

\textbf{Conclusion}: Inclusion measures are grounded in the real barriers identified
for this academic context --- knowledge scaffolding and typographic accessibility ---
while also addressing the broader gender gap through ethical framing and accessible
entry points.

\medskip
\textbf{Overall Feasibility.} The project is feasible across all evaluated dimensions. Technical risks have been
mitigated through hardware investment (64~GB RAM, additional module acquired during
development). The budget is within academic scope and the cost model is transparent.
Legal and ethical obligations are satisfied through open-source licensing and a locally
controlled infrastructure. All planned activities fit within the TecnoCampus normative
timeline. No identified risk is considered project-blocking.

Full phase task lists (Task IDs, hours, critical flags), the five-phase NIST SP 800-30
risk assessment, and detailed budget breakdowns are provided in
the \href{https://theegea.github.io/TFG/annexos/app-i/}{project online documentation}.
